\section{Verlinde 公式：融合规则为什么被模变换决定}
\label{sec:verlinde}

\subsection{这一章要解决什么}

二维有理共形场论（rational conformal field theory，RCFT）给每个可积表示 $R_i$ 配一条融合规则
\begin{equation}
R_i\otimes R_j=\bigoplus_k N_{ij}{}^{k}R_k ,
\label{eq:fusion-v}
\end{equation}
其中重数 $N_{ij}{}^{k}$ 是整数。另一面，环面 $T^2$ 上的模变换生成元 $S$、$T$ 作用在真空模空间上，给出一个矩阵 $S_{ij}$。Verlinde 在 1988 年提出：这两组数据不是独立的，$N_{ijk}$ 完全由 $S$ 决定，
\begin{equation}
N_{ijk}=\sum_l\frac{S_{il}S_{jl}S_{kl}}{S_{0l}} .
\label{eq:verlinde}
\end{equation}
本章要回答四个问题。第一，Witten 在 \S4.5 是怎样从三维 Chern--Simons 理论“推出”这个公式的？他的推导可以拆成十步，每一步都只用到非常少的输入，本章要把这十步写成可逐条核对的式子。第二，这条推导到底用了什么、没用？一个流行的误解是它用了共形块满足 Knizhnik--Zamolodchikov 方程，事实正好相反：整条链没有一步用到共形块的动力学。第三，(4.51) 有两种读法，一种把 $S$ 当作派生量，一种把 $S$ 当作基本量；两种读法数学等价，但解释不同，本章要把它们分开。第四，“Witten 证明了 Verlinde 猜想”这句话不准确：本文做的是物理推导，数学证明是另一条线（TUY 1989、Faltings 1994），本章最后把四条路线的严格性排成一张表。

\subsection{原文说了什么}

\paragraph{（一）把待求的量翻译成三维 Hilbert 空间的维数。}
原文先在 $S^2\times S^1$ 上工作。设 $\langle R\rangle$ 是沿 $\{P_a\}\times S^1$ 的一组平行 Wilson 线，表示分别为 $R_a$。把 $S^1$ 方向的演化看成迹，路径积分等于 $S^2$ 上带电荷的 Hilbert 空间的维数：
\begin{equation}
Z(S^2\times S^1;\langle R\rangle)=\dim\cH_{S^2;\langle R\rangle}. \tag{4.33}
\end{equation}
于是单个电荷、两个电荷、三个电荷分别给出
\begin{equation}
Z(S^2\times S^1;R_a)=\delta_{a,0}, \tag{4.34}
\end{equation}
\begin{equation}
Z(S^2\times S^1;R_a,R_b)=g_{ab}, \tag{4.35}
\end{equation}
\begin{equation}
Z(S^2\times S^1;R_a,R_b,R_c)=N_{abc}. \tag{4.36}
\end{equation}
原文在 (4.36) 后紧接着说：\emph{“with Nabc the trilinear `coupling' of Verlinde, since this is the dimension of the physical Hilbert space.”} \src{388} 这句话是整章的关键：$N_{abc}$ 在这里不是从融合规则拿来的，而是被\emph{定义}成三维物理 Hilbert 空间的维数。后面要证的，是这个三维定义的量与 $S$ 矩阵的三重积相等。$g_{ab}$ 是电荷共轭度规：$R_b$ 是 $R_a$ 的对偶表示时为 $1$，否则为 $0$；它负责升降指标，使 $N_{abc}=N_{ab}{}^{c}$ 这类写法有意义。

\paragraph{（二）手术把 $S^3$ 与 $S^2\times S^1$ 连起来。}
沿 $S^2\times S^1$ 中一个圈 $C$ 的管状邻域（环面）做 Dehn 手术。取模变换
\begin{equation}
S:\quad a\mapsto b,\qquad b\mapsto -a
\end{equation}
（在基 $a,b$ 下矩阵为 $\begin{psmallmatrix}0&-1\\1&0\end{psmallmatrix}$，等价于 $\tau\mapsto-1/\tau$ \src{389}），手术后的三维流形正是 $S^3$。用 (4.27)（沿 $C$ 做手术时配分函数按 $S$ 矩阵展开）：
\begin{equation}
Z(S^3)=\sum_j S_0{}^j\,Z(S^2\times S^1;R_j). \tag{4.37}
\end{equation}
代入 (4.34) 得
\begin{equation}
Z(S^3)=S_{0,0}. \tag{4.38}
\end{equation}
对 $SU(2)_k$，$S$ 有显式公式（$m,n=0,1,\dots,k$ 对应自旋 $m/2,n/2$）：
\begin{equation}
S_{mn}=\sqrt{\frac{2}{k+2}}\,\sin\frac{\pi(m+1)(n+1)}{k+2}.
\label{eq:smatrix}
\end{equation}
原文对一般 $G$ 只给出 $S$ 由 $q=\exp\frac{2\pi\ii}{k+h^\vee}$ 决定，$SU(N)$ 时 $q=\exp\frac{2\pi\ii}{N+k}$ \src{382}；上式是 $SU(2)$ 的特例 \src{390}。

\paragraph{（三）给 $S^3$ 加 Wilson 线。}
在 $S^2\times S^1$ 上让一条 $R_j$ 线平行于手术圈 $C$，手术后它在 $S^3$ 上仍是 $R_j$ 线，而手术圈变成与它相链的圈。用 (4.27)：
\begin{equation}
Z(S^3;R_j)=\sum_i S_0{}^i\,Z(S^2\times S^1;R_i,R_j). \tag{4.42}
\end{equation}
再用 (4.35)：
\begin{equation}
Z(S^3;R_j)=\sum_i S_0{}^i g_{ij}=S_{0,j}. \tag{4.43}
\end{equation}
取 $SU(2)$、$R$ 为二维表示，与无结 Wilson 线的期望值 (4.23) 比较，得到
\begin{equation}
\langle C\rangle=\frac{S_{0,1}}{S_{0,0}}
=\frac{\sin\frac{2\pi}{k+2}}{\sin\frac{\pi}{k+2}}. \tag{4.44}
\end{equation}
两条互不链接的线则给
\begin{equation}
Z(S^3;R_j,R_k)=\sum_i S_0{}^i N_{ijk}, \tag{4.45}
\end{equation}
\begin{equation}
Z(S^3;R_j,R_k)=\frac{S_{0,j}S_{0,k}}{S_{0,0}}. \tag{4.46}
\end{equation}
注意 (4.46) 已经是 (4.51) 的一个特例：把 (4.51) 的指标 $i$ 取成 $0$ 就得到它。

\paragraph{（四）斜体小标题之后的推导。}
原文在 p.391 有一段斜体小标题：
\begin{quote}
\emph{Proof of Verlinde's Conjecture.} The last equation is a special case of a celebrated conjecture by Verlinde, which has been proved by Moore and Seiberg [13]. We can use these methods to give a new proof of Verlinde's conjecture, in the case of current algebra.
\end{quote}
\src{391} 这段引文要逐字读：猜想\emph{已由} Moore--Seiberg 证明，Witten 给自己的定位是“在流代数情形给出\emph{新证明}”。同页还说明将使用“广义手术关系 (4.28)”。这个公式早在 p.386 给出时就预告过：\emph{“This formula will be used later in a new proof of Verlinde's conjecture.”} \src{386} 接着原文在 $S^3$ 上放一对相链的 Wilson 线 $L(R_i;R_j)$：$R_i$ 沿手术圈，$R_j$ 沿与它平行的圈。由 (4.28)，
\begin{equation}
Z(S^3;L(R_i;R_j))=\sum_k S_i{}^k\,Z(S^2\times S^1;R_k,R_j)=S_{ij}. \tag{4.47}
\end{equation}
这里 $S_i{}^k g_{kj}=S_{ij}$ 用了 $g$ 降指标与 $S$ 的对称性。若再把 $R_i$ 线与另外两条线 $R_j,R_k$ 相链，则
\begin{equation}
Z(S^3;L(R_i;R_j,R_k))=\sum_m S_i{}^m N_{mjk}. \tag{4.48}
\end{equation}
关键的几何观察是：$L(R_i;R_j,R_k)$ 可以看成 $L(R_i;R_j)$ 与 $L(R_i;R_k)$ 沿 $R_i$ 线做连通和。于是
\begin{equation}
Z(S^3;L(R_i;R_j,R_k))\,Z(S^3;R_i)
=Z(S^3;L(R_i;R_j))\,Z(S^3;L(R_i;R_k)). \tag{4.49}
\end{equation}
原文把这一步的依据称为“一维线性代数”：$S^2$ 上带两个 $R_i$ 电荷的 Hilbert 空间是一维的，一维空间里的任何配对都自动分解为乘积。把 (4.47)(4.48)(4.43)(4.38) 代入 (4.49)：
\begin{equation}
Z(S^3;L(R_i;R_j,R_k))=\frac{S_{ij}S_{ik}}{S_{0,i}}, \tag{4.50}
\end{equation}
与 (4.48) 联立即得
\begin{equation}
\sum_m S_i{}^m N_{mjk}=\frac{S_{ij}S_{ik}}{S_{0,i}} .
\label{eq:verlinde-key}
\tag{4.51}
\end{equation}

\paragraph{（五）幂等元。}
原文没有停在 (4.51)，而是把它改写成融合代数的一组幂等元。令 $v_i$ 是 Verlinde 基（$v_0=\chi$ 为真空），融合代数为
\begin{equation}
v_i v_j=\sum_k N_{ij}{}^{k}v_k ,
\end{equation}
定义
\begin{equation}
w_i=S_{0,i}\sum_m S_i{}^m v_m . \tag{4.52}
\end{equation}
原文断言
\begin{equation}
w_i w_j=\delta_{ij}w_j , \tag{4.54}
\end{equation}
中间一步 (4.53) 在印本中写作 $w_iw_j=S_j{}^l v^m\cdot S_{il}S_i{}^m\cdot S_{0,j}$，其中 $v^m$ 应为 $v_m$（这一处印本可疑，见 G9.2）。这就是“两种读法”的第二种：不只是说 $S$ 给出 $N$，而是说融合代数被 $w_i$ 对角化。

\subsection{原文没说的细节}

\begin{enumerate}[label=G9.\arabic*,leftmargin=*]
  \item \textbf{十步链没有被写成步骤。}原文按几何操作顺序叙述（先算 $S^3$，再加一条线、两条线、三条线），每一步只说一句“应用 (4.27)”“应用 (4.28)”。全链真正的新输入只有三处：手术律 (4.27)/(4.28)、$N_{abc}$ 的\emph{定义} (4.36)、以及一维 Hilbert 空间的秩一性 (4.49)。本章表 \ref{tab:tenstep} 把它们排成十步，并标出每步用的是哪一处输入。\ours
  \item \textbf{(4.51) 与三重积的等价性没有写。}(4.51) 是“一个指标在上”的形式，\eqref{eq:verlinde} 是完全对称的三重积形式。两者等价要用 $S$ 的幺正性 $\sum_i S_i{}^m S_i{}^n=\delta^{mn}$，原文没有给出这一步；本章 D-9.3 补出双向推导。另外 (4.52)$\to$(4.54) 之间印本跳过太多，且 (4.53) 有 $v^m/v_m$ 的记号问题。\ours
  \item \textbf{四种“证明”的地位没有区分。}本文只引 Moore--Seiberg [13]，并把 TUY 1989 与 Faltings 1994 留给数学文献；Reshetikhin--Turaev 1991 更在本文之后。原文没有说“物理推导”与“数学证明”差在哪里，只用了 “proof” 一词。本章表 \ref{tab:routes} 按工具与严格性把四条路线分开。\ours
  \item \textbf{$SU(2)$ 只给了 $S$ 的公式，没有给 $N$ 的表。}原文用 (4.39) 算出了 (4.44) 的比值，但没有把 $N_{ijk}$ 列成表，也没有讨论 $k=1,2$ 的截断后果。本章表 \ref{tab:su2} 补出 $k=1,2$ 的完整复算。\ours
  \item \textbf{截断机制没有解释。}可积表示的条件是 $2j\le k$：这是仿射 $\widehat{\mathfrak{sl}}_2$ 表示可积性的要求，不是融合规则本身的选择。求和指标 $m$ 只跑可积表示，所以 $N_{ijk}$ 在 $2j>k$ 时为零。$N_{\frac12,\frac12,\frac12}=0$（$k=1$）有两个层次的原因，容易混：普通 $SU(2)$ 里 $\frac12\otimes\frac12$ 含自旋 $1$，但 $k=1$ 时 $2\cdot1>k$，自旋 $1$ 不可积，于是被切掉；到了 $k=2$，同一个 $N_{\frac12,\frac12,1}$ 就等于 $1$。\ours
\end{enumerate}

\subsection{补出的推导}

先把十步链排成表。第 $0$ 步是出发点，第 $1$–$10$ 步是原文的推导顺序。

\begin{table}[htbp]
\centering
\small
\caption{Verlinde 公式的十步链（第 $0$ 步为出发点）}
\label{tab:tenstep}
\begin{tabular}{@{}c >{\raggedright\arraybackslash}p{8.1cm} c >{\raggedright\arraybackslash}p{4.3cm}@{}}
\toprule
步 & 操作／式子 & 页 & 这一步靠什么 \\
\midrule
0 & $Z(X\times S^1)=\Tr_{\cH_X}(\mathbf 1)=\dim\cH_X$ \ (4.30) & 387 & 三维理论哈密顿量为零，演化算符是恒等 \\
1 & $Z(\widetilde M)=\sum_j K_0{}^j Z(M;R_j)$ \ (4.27) & 385 & 沿手术圈 $C$ 插入 $R_j$ 线（手术律）\\
2 & $Z(\widetilde M;R_i)=\sum_j K_i{}^j Z(M;R_j)$ \ (4.28) & 386 & 同上，且 $R_i$ 线不随手术消失 \\
3 & $S^3=(S^2\times S^1)$ 沿 $C$ 做 $S$ 手术；$S^2\times S^1$ 由平移手术得到（图 15a、16） & 389 & Heegaard 分解／Dehn 手术 \\
4 & 单电荷 $\delta_{a0}$、双电荷 $g_{ab}$、三电荷 $N_{abc}$ \ (4.34)--(4.36) & 388 & $N$ 的\emph{定义}：物理 Hilbert 空间维数 \\
5 & $Z(S^3)=\sum_j S_0{}^j Z(S^2\times S^1;R_j)=S_{0,0}$ \ (4.37)(4.38) & 389 & 第 1 步 + 第 4 步（代 (4.34)）\\
6 & $Z(S^3;R_j)=S_{0,j}$ \ (4.42)(4.43) & 391 & 第 2 步 + 第 4 步（代 (4.35)）\\
7 & $Z(S^3;L(R_i;R_j))=\sum_k S_i{}^k g_{kj}=S_{ij}$ \ (4.47) & 392 & 第 2 步；线随手术保留成链接对 \\
8 & $Z(S^3;L(R_i;R_j,R_k))=\sum_m S_i{}^m N_{mjk}$ \ (4.48) & 392 & 第 2 步 + 第 4 步 \\
9 & (4.49)$\Rightarrow$(4.50)：$Z(L(R_i;R_j,R_k))=S_{ij}S_{ik}/S_{0,i}$ & 393 & 连通和 + 一维 Hilbert 空间秩一 \\
10 & (4.51)；幂等元 (4.52)--(4.54) & 393 & 第 9 步 + $S$ 正交性 \\
\bottomrule
\end{tabular}
\end{table}

\begin{derivation}[D-9.1：第 7 步，(4.47) 的来历]
在 $S^2\times S^1$ 上取一条 $R_j$ 线平行于手术圈 $C$，并让 $C$ 上的线标记为 $R_i$（图 15b）。对 $C$ 做手术，$S^2\times S^1$ 变成 $S^3$，原来的 $R_j$ 平行线与新形成的 $R_i$ 圈相链。按 (4.28)，手术前应对称地写出
\[
Z(S^3;L(R_i;R_j))=\sum_k S_i{}^k\,Z(S^2\times S^1;R_k,R_j).
\]
其中 $S_i{}^k$ 的\emph{上}指标 $k$ 来自手术圈上的表示，\emph{下}指标 $i$ 来自 $S^3$ 中的 $R_i$ 线；这是 (4.28) 中 $K_i{}^j$ 的指标约定。由 (4.35)，
\[
Z(S^2\times S^1;R_k,R_j)=g_{kj}.
\]
于是
\[
Z(S^3;L(R_i;R_j))=\sum_k S_i{}^k g_{kj}.
\]
$g$ 的作用是把上指标降下来：定义 $S_{ij}:=\sum_k S_i{}^k g_{kj}$。对 $SU(2)$，$g_{ij}=\delta_{i\bar j}$，且所有表示自对偶，故 $S_{ij}=S_i{}^j$，并且 $S$ 对称。对一般 $G$，$S_{ij}$ 对称而 $S_i{}^j$ 不必然对称，二者不能混写。这一步是全链唯一需要 (4.28)——而不仅是 (4.27)——的地方：平行线给 $g$，链接线给 $S$，手术把“平行”换成“链接”。\ours
\end{derivation}

\begin{derivation}[D-9.2：第 9 步，秩一性 (4.49)$\Rightarrow$(4.50)]
先说明为什么 (4.49) 成立。$S^2$ 上带两个电荷 $R_i,\bar R_i$ 的 Hilbert 空间是一维的：由 (4.35)，$Z(S^2\times S^1;R_i,\bar R_i)=g_{i\bar i}=1$，而 $\dim\cH=Z$，所以 $\cH$ 一维。
取 $\cH$ 的一组基 $e$ 与对偶基 $e^*$（$\langle e^*,e\rangle=1$）。一维空间中的向量与对偶向量都可以写成
\[
\psi=\alpha e,\quad \chi=\beta e^*,\qquad v'=\gamma e^*,\quad v=\delta e .
\]
四个配对分别是
\[
(\psi,\chi)=\alpha\beta,\quad (v',v)=\gamma\delta,\qquad
(\psi,v)=\alpha\delta,\quad (v',\chi)=\gamma\beta .
\]
于是
\[
(\psi,\chi)(v',v)=\alpha\beta\gamma\delta=(\psi,v)(v',\chi),
\]
这正是 (4.49)。几何上，$L(R_i;R_j,R_k)$ 沿 $R_i$ 线剪开成 $L(R_i;R_j)$ 与 $L(R_i;R_k)$ 的连通和，剪开处插入的完备基就是 $\sum_n|e_n\rangle\langle e_n|$；一维时只剩一项，所以四个配对按上式分解。
再算 (4.50)：把 (4.47)(4.48)(4.43)(4.38) 代入 (4.49)，
\[
\Big(\sum_m S_i{}^m N_{mjk}\Big)\cdot S_{0,i}=S_{ij}\cdot S_{ik}.
\]
两边除以 $S_{0,i}$（$SU(2)$ 时 $S_{0,i}\neq0$）即得 (4.50)。\ours
\end{derivation}

\begin{derivation}[D-9.3：(4.51) 与三重积形式等价]
\textbf{从 (4.51) 到三重积。}把 (4.51) 两边乘以 $S_i{}^n$ 并对可积指标 $i$ 求和：
\[
\sum_i S_i{}^n\sum_m S_i{}^m N_{mjk}=\sum_i S_i{}^n\frac{S_{ij}S_{ik}}{S_{0,i}} .
\]
左边用 $S$ 的正交性 $\sum_i S_i{}^n S_i{}^m=\delta^{nm}$，得 $\sum_m\delta^{nm}N_{mjk}=N_{njk}$。右边把 $S_i{}^n$ 写成 $S_{ni}$（$SU(2)$；一般 $G$ 用 $S_{ij}$ 的对称性），再把求和指标 $i$ 改名为 $l$：
\[
N_{njk}=\sum_l\frac{S_{nl}S_{lj}S_{lk}}{S_{0,l}},
\]
即 \eqref{eq:verlinde}（把 $n$ 改写为 $i$）。

\textbf{从三重积到 (4.51)。}把 $N_{mjk}=\sum_l S_{ml}S_{jl}S_{kl}/S_{0,l}$ 代入 (4.51) 左边：
\[
\sum_m S_i{}^m N_{mjk}
=\sum_l\frac{S_{jl}S_{kl}}{S_{0,l}}\sum_m S_i{}^m S_{ml}
=\sum_l\frac{S_{jl}S_{kl}}{S_{0,l}}\,\delta_{il}
=\frac{S_{ji}S_{ki}}{S_{0,i}} .
\]
这正是 (4.51) 的右边。所以两种写法等价，但 (4.51) 的指标位置更适合接手术链，\eqref{eq:verlinde} 的三重积形式更适合做代数计算。\ours
\end{derivation}

\begin{derivation}[D-9.4：(4.52)$\Rightarrow$(4.54) 的中间步]
在 $SU(2)$（所有表示自对偶，$g_{ij}=\delta_{ij}$，故 $N_{ijk}=N_{ij}{}^k$ 全对称）中计算；一般 $G$ 把 $N$ 换成全对称的三指标版本，步骤相同。把 $v_mv_n=\sum_k N_{mn}{}^k v_k$ 代入 (4.52) 的乘积：
\[
w_iw_j=S_{0,i}S_{0,j}\sum_{m,n}S_i{}^m S_j{}^n\,v_mv_n
=S_{0,i}S_{0,j}\sum_{m,n,k}S_i{}^m S_j{}^n N_{mn}{}^k\,v_k .
\]
用全对称性 $N_{mn}{}^k=N_{mnk}$，并对 $n$ 用 (4.51) 的等价形式。把 (4.51) 中的指标重命名：$i\to j,\ m\to n,\ j\to m$，得
\[
\sum_n S_j{}^n N_{nmk}=\frac{S_{jm}S_{jk}}{S_{0,j}} .
\]
（这一步用到了 $N_{nmk}=N_{mnk}$；若不做全对称假设，就要先把 $N_{mn}{}^k$ 用 $g$ 降成 $N_{mnk}$，再交换指标。）代回：
\[
w_iw_j=S_{0,i}S_{0,j}\sum_{m,k}S_i{}^m\frac{S_{jm}S_{jk}}{S_{0,j}}\,v_k
=S_{0,i}\Big(\sum_m S_{im}S_{jm}\Big)\sum_k S_{jk}v_k .
\]
由 $S$ 的正交性 $\sum_m S_{im}S_{jm}=\delta_{ij}$，
\[
w_iw_j=\delta_{ij}S_{0,i}\sum_k S_{jk}v_k=\delta_{ij}w_j .
\]
最后一步用了 $i=j$ 时 $S_{0,i}=S_{0,j}$。这就是 (4.54)。印本 (4.53) 把中间式写成 $v^m$ 是笔误：整个推导用的是 $v_m$，且 $S_{il}S_i{}^m$ 的指标配对也只有在 $v_m$ 下才自洽。\ours
\end{derivation}

\begin{derivation}[D-9.5：$SU(2)$ 的显式复算]
用 (4.39) 与 \eqref{eq:verlinde} 对 $k=1,2,3$ 做数值复算（脚本 \texttt{\_w2\_check\_su2.py}，双精度，正交性误差 $<10^{-16}$）。

\textbf{$k=1$}：$j\in\{0,\frac12\}$，
\[
S=\frac{1}{\sqrt2}\begin{pmatrix}1&1\\1&-1\end{pmatrix}.
\]
非零融合系数只有
\[
N_{000}=1,\quad N_{0\frac12\frac12}=1,\quad N_{\frac12 0\frac12}=1,\quad N_{\frac12\frac12 0}=1,
\]
而 $N_{\frac12\frac12\frac12}=0$。后一个零可直接验证：
\[
N_{\frac12\frac12\frac12}
=\sum_{l=0,\frac12}\frac{S_{\frac12 l}^3}{S_{0l}}
=\frac{(1/\sqrt2)^3}{1/\sqrt2}+\frac{(-1/\sqrt2)^3}{1/\sqrt2}
=\frac12-\frac12=0 .
\]

\textbf{$k=2$}：$j\in\{0,\frac12,1\}$，
\[
S=\begin{pmatrix}
\frac12 & \frac{1}{\sqrt2} & \frac12\\[2pt]
\frac{1}{\sqrt2} & 0 & -\frac{1}{\sqrt2}\\[2pt]
\frac12 & -\frac{1}{\sqrt2} & \frac12
\end{pmatrix}.
\]
按 \eqref{eq:verlinde} 算出的非零融合系数恰有 $10$ 个，全部等于 $1$：
\[
(0,0,0),\ (0,1,1),\ (0,2,2),\ (1,0,1),\ (1,1,0),\ (1,1,2),\ (1,2,1),\ (2,0,2),\ (2,1,1),\ (2,2,0),
\]
其中指标 $0,1,2$ 分别代表自旋 $0,\frac12,1$。特别地 $N_{111}=0$、$N_{221}=0$、$N_{222}=0$：$\frac12\otimes\frac12=0\oplus1$ 不含 $\frac12$；而 $1\otimes1$ 在能级 $2$ 只含 $0$——$j=1,2$ 都被截断条件 $j\le k-j_1-j_2=0$ 切掉。这与截断融合规则
\[
j_1\otimes j_2=\bigoplus_{j=|j_1-j_2|}^{\min(j_1+j_2,\,k-j_1-j_2)}j
\]
逐条一致。

\textbf{$k=3$}：可积表示 $j=0,\frac12,1,\frac32$。复算给出 $N_{112}=N_{123}=N_{222}=1$，而 $N_{133}=0$（$\frac12\otimes\frac32=1\oplus2$ 在能级 $3$ 只含 $1$）。这三组数与手算的截断规则完全一致。\ours
\end{derivation}

\begin{table}[htbp]
\centering
\small
\caption{$SU(2)_k$ 的 $S$ 矩阵与非零融合系数（复算结果）}
\label{tab:su2}
\begin{tabular}{@{}c >{\raggedright\arraybackslash}p{5.2cm} >{\raggedright\arraybackslash}p{7.6cm}@{}}
\toprule
$k$ & $S$ 矩阵 & 非零 $N_{ijk}$（指标 $m$ 代表自旋 $m/2$）\\
\midrule
$1$ & $\dfrac{1}{\sqrt2}\begin{pmatrix}1&1\\1&-1\end{pmatrix}$
& $N_{000}=N_{011}=N_{101}=N_{110}=1$；$N_{111}=0$ \\
$2$ & $\begin{pmatrix}\frac12&\frac{1}{\sqrt2}&\frac12\\
\frac{1}{\sqrt2}&0&-\frac{1}{\sqrt2}\\
\frac12&-\frac{1}{\sqrt2}&\frac12\end{pmatrix}$
& $N_{000},N_{011},N_{022},N_{101},N_{110},$\newline
$N_{112},N_{121},N_{202},N_{211},N_{220}=1$；其余为 $0$ \\
\bottomrule
\end{tabular}
\end{table}

\subsection{与其他来源的对照}

\begin{center}
\begin{minipage}{\linewidth}
\centering
\captionof{table}{四条路线的严格性}
\label{tab:routes}
\small
\small
\setlength{\tabcolsep}{4pt}
\begin{tabular}{@{}p{2.6cm} p{3.1cm} p{1.9cm} p{4.3cm} p{2.7cm}@{}}
\toprule
层级 & 代表 & 年代 & 工具 & 严格性 \\
\midrule
物理推导 & Witten \S4.5（本文） & 1988/89 & TQFT 公理、手术、$N$ 的三维定义 & 物理推导，可复算 \\
物理证明 & Moore--Seiberg & 1989 & 模不变性、融合矩阵自洽性 & 物理证明（本文承认）\\
数学证明 & TUY；Faltings & 1989；1994 & 共形块丛 + KZ 射影平坦联络；rank--level 对偶 & 数学定理 \\
组合／范畴 & Reshetikhin--Turaev 及后续 & 1991 起 & 量子群 $R$ 矩阵、模张量范畴 & 独立构造，$S$ 为公理数据 \\
\bottomrule
\end{tabular}
\end{minipage}
\end{center}

这四条路线的差别不是“谁对谁错”，而是\emph{证明了什么}。Witten 的推导从三维理论出发，用配分函数在手术下的变换性质；它证明的是“在 CS 理论自洽的前提下，(4.51) 必须成立”。Moore--Seiberg 的证明从二维 RCFT 的模不变性出发，证明的是“RCFT 的融合矩阵满足 (4.51)”。TUY 与 Faltings 从代数几何出发，把共形块丛的维数与 Verlinde 公式的右端相等。RT 则把 $S$ 当作模张量范畴的公理数据，Verlinde 公式成为范畴的定理。四者结论相同，前提不同。

\begin{center}
\begin{minipage}{\linewidth}
\centering
\captionof{table}{同一公式在不同来源中的位置}
\label{tab:sources}
\small
\small
\setlength{\tabcolsep}{4pt}
\begin{tabular}{@{}>{\raggedright\arraybackslash}p{2.2cm} >{\raggedright\arraybackslash}p{3.3cm} >{\raggedright\arraybackslash}p{2.7cm} >{\raggedright\arraybackslash}p{2.6cm} >{\raggedright\arraybackslash}p{2.7cm}@{}}
\toprule
主题 & 本文 & 河野 & 上野 & Schottenloher \\
\midrule
公式形式 & (4.51)、(4.52)--(4.54) \src{393} & 命题 1.22，书 p.37 & 例 5.23，p.106 & 定理 11.6，p.177 \\
$S$ 矩阵 & (4.39) \src{390} & 书 p.60 & (5.29)(5.30)，p.103 & p.177 \\
证明路线 & 手术推导 \src{385--393} & 定理 2.16 以 1.22 为砖，p.60 & 引 Moore--Seiberg 首证，p.106 & TUY$\to$Faltings$\to$组合，p.171--178 \\
幂等元 & (4.52)--(4.54) & 命题 1.22 同一构造 & 未核 & \S11.3 谱分解 \\
截断 & $2j\le k$ \src{382} & 书 p.60 & p.106 & p.177 \\
\bottomrule
\end{tabular}
\end{minipage}
\end{center}

河野书里同一个 $S$ 矩阵出现两次，角色不同：命题 1.22（书 p.37）是融合代数的基变换，就是 Verlinde 公式；定理 2.16（书 p.60）是 Dehn 手术核，用来定义 Witten 不变量。定理 2.16 的证明拿 1.22 当砖，但把手术公式当成“Verlinde 的证明”是张冠李戴。上野书把 $S$ 写成 $S_{\lambda\mu}=S_{\lambda0}\,\mathrm{ch}_\mu\big(e^{-2\pi\ii(\lambda+\rho)/(\ell+g^*)}\big)$，$S_{\lambda0}$ 是 $\sin$ 的乘积（(5.29)(5.30)，p.103），并明确说 Verlinde 猜想“最先由 Moore--Seiberg 证明”（定理 5.22 前，p.106）。Schottenloher 把公式写成配分函数的 $\sin$ 幂次和（定理 11.6，p.177），并说这些证明“都不容易理解”（p.171）：KZ 单值性加非交换 Hodge 理论是重机器，而右端是整数这件事本身也不显然。本讲未核李思书中的对应表述（见本章边界）。

\subsection{陷阱与约定}

\begin{keypoint}
\textbf{(4.36) 是定义，(4.51) 是定理，(4.46) 只是特例。}读原文时最容易把这三者的逻辑顺序搞反：$N_{abc}$ 先被定义成三维 Hilbert 空间的维数（(4.36)，与 $S$ 无关），然后十步链证明它等于 $S$ 的三重积（(4.51)）；(4.46) 只是 (4.51) 取 $i=0$ 的特例，不能当成“证明”。另外，“Proof of Verlinde's Conjecture” 是 p.391 的\emph{段首斜体小标题}，不是节标题；\S4.5 的节标题是 “Some Concrete Surgeries”（p.388）。标题本身就说明这一节的主线是手术演算，Verlinde 公式是副产品。
\end{keypoint}

\begin{itemize}[leftmargin=1.6em]
  \item \textbf{指标位置。}$S_i{}^j$（一上一下）与 $S_{ij}$（两下）在 $SU(2)$ 相等，在一般 $G$ 不等：$S_{ij}=\sum_k S_i{}^k g_{kj}$。$N_{ij}{}^k$ 与 $N_{ijk}$ 也差一个 $g$。原文 (4.51) 左边是 $S_i{}^mN_{mjk}$，三个指标全在下，不能随手把 $S_i{}^m$ 写成 $S_{im}$ 而不说明 $g$ 的约定。
  \item \textbf{归一化。}$Z(S^3)=S_{0,0}$，不是 $1$。$SU(2)$ 时 $S_{0,0}=\sqrt{2/(k+2)}\sin\frac{\pi}{k+2}$；$k=2$ 时它是 $\frac12$。所有“比值”公式（如 (4.44)）依赖这个归一化，单独写 $Z(S^3;R_j)=S_{0,j}$ 时也一样。
  \item \textbf{$q$ 的定义。}$SU(N)$ 用 $q=\exp\frac{2\pi\ii}{N+k}$ \src{382}，一般 $G$ 用 $q=\exp\frac{2\pi\ii}{k+h^\vee}$，$SU(2)$ 时二者都是 $\exp\frac{2\pi\ii}{k+2}$。不要写成 $q=\exp\frac{2\pi\ii}{k}$。
  \item \textbf{标架反常。}原文 (4.41) 给出标架变换下的相位 $e^{2\pi\ii(n-m)c/24}$，并明说在本文的简单应用中“没有必要”追踪 \src{390--391}。全链用的是标准手术 $S$；若要换成 $T^nST^m$，(4.51) 右边会多相位，但 (4.50) 的比值里 $c/24$ 相位相消。本讲不处理一般标架。
  \item \textbf{“秩一性”的准确含义。}不是“某个 Hilbert 空间维数为 $1$”，而是 $\dim\cH_{S^2;R_i,\bar R_i}=1$；等价地，(4.35) 的 $g_{i\bar i}=1$。若把“维数一”错误地安到三电荷空间上，就会把 $N_{ijk}$ 当成 $0$ 或 $1$ 的平凡量。
  \item \textbf{截断与指标范围。}$m$ 只跑 $0,\dots,k$（$SU(2)$）。求和指标和结果指标都要截断；只截一边会得到非整数。$N_{111}=0$（$k=1$）与 $N_{221}=0$（$k=2$）都是截断，不是“融合规则不含该表示”的普遍事实。
\end{itemize}

\subsection{练习}

\begin{enumerate}[leftmargin=*]
  \item （复算）用 (4.39) 写出 $k=2$ 的 $S$，验证 $S^{\mathsf T}=S$ 与 $S^2=\mathbf 1$。再用 \eqref{eq:verlinde} 算 $N_{\frac12\frac12 0}$、$N_{\frac12\frac12 1}$、$N_{111}$，与表 \ref{tab:su2} 对照。
  \item （复算）取 $k=3$，验证 $N_{112}=1$、$N_{123}=1$、$N_{222}=1$、$N_{133}=0$。用截断融合规则 $j_1\otimes j_2=\bigoplus_{j=|j_1-j_2|}^{\min(j_1+j_2,k-j_1-j_2)}j$ 解释这四个数。
  \item （推导）从 (4.51) 取 $i=0$，推出 (4.46)。说明这一步为什么不能反过来当成 (4.51) 的证明。
  \item （推导）补出 D-9.3 的两个方向，并指出哪一步用了 $S$ 的幺正性、哪一步用了 $S$ 的对称性。
  \item （推导）补出 (4.52)$\to$(4.54) 的全部中间步，并说明印本 (4.53) 的 $v^m$ 为什么必须改成 $v_m$。
  \item （判断）“Witten 在 \S4.5 证明了 Verlinde 猜想。”这句话对吗？请分别就“物理推导”“数学证明”两个层面说明，并指出原文 p.391 对 Moore--Seiberg 的表述。
  \item （判断）“这条十步链用到了共形块满足 KZ 方程。”对吗？请列出链中真正用到的三处输入，并说明为什么共形块动力学没有进入。
  \item （判断）“$N_{\frac12\frac12\frac12}=0$ 说明 $\frac12\otimes\frac12$ 不含自旋 $1$。”这个推理在 $k=1$ 与 $k=2$ 下分别对不对？请用截断条件 $2j\le k$ 说明。
\end{enumerate}

\begin{boundary}
\begin{itemize}[leftmargin=1.4em]
  \item \textbf{原文证明：}$N_{abc}$ 的三维定义 (4.36)；手术链 (4.37)--(4.51)；(4.52)--(4.54) 的幂等元陈述；$SU(2)$ 的 $S$ 公式 (4.39)。原文的全部应用只用标准手术 $S$，标架反常 (4.41) 被明确搁置 \src{390--391}。
  \item \textbf{我们补的：}表 \ref{tab:tenstep} 的十步分解与三处输入的归纳；D-9.1--D-9.5 的完整中间步；两种读法的方向区分；$SU(2)_{k=1,2,3}$ 的数值复算；河野命题 1.22 与定理 2.16 的角色区分；四条路线的严格性表。
  \item \textbf{未证／未核：}数学严格证明（TUY 1989、Faltings 1994）只按 Schottenloher 的转述（p.171--178）与上野 p.106 的引述，未读原文；河野 p.37、p.60 的公式未逐字核对；李思书中的对应表述未核；印本 (4.53) 的 $v^m$ 判为笔误，但未见勘误表；一般 $G$ 下 $N_{ijk}$ 的全对称性与 $g$ 的约定只按 D-9.1 的说明，未逐群验证。
  \item \textbf{练习答案：}本讲只给提示，未写入附录 C（本讲的写权限不含 \texttt{appC}）。
\end{itemize}
\end{boundary}
